% =====================================================================================
% math_formulation.tex
% Standalone, publication-ready derivations for MultiMax activation probing
% and uncertainty-gated cascades.
%
% Requires: amsmath, amssymb, amsthm, mathtools.
% Every environment is self-contained; paste sections directly into a paper body.
% Cross-references use \label{...} keys prefixed eq:/thm:/lem:/prop:/asm:.
% =====================================================================================

\section{Token-Level Feature Extraction and Probing Pipelines}
\label{sec:pipeline}

Let $M$ denote a frozen autoregressive transformer and let $x_{i,j}\in\mathbb{R}^{d}$ be
the residual-stream activation at layer $\ell$, sequence position $j\in\{1,\dots,N_i\}$,
for input $i$. A \emph{probe} is a map $f_\theta:\mathbb{R}^{N_i\times d}\to\mathbb{R}$
producing a scalar logit, composed of a token-level transform $\phi$ and a permutation
sequence-reduction $\rho$:
%
\begin{equation}
  f_\theta(X_i)\;=\;\rho\!\big(\phi(x_{i,1}),\,\dots,\,\phi(x_{i,N_i})\big),
  \qquad
  \phi:\mathbb{R}^{d}\to\mathbb{R}^{m}.
  \label{eq:probe-general}
\end{equation}
%
All probes considered here share $\phi$ and differ only in $\rho$. This is the point of
the construction: it isolates the \emph{aggregation rule} as the independent variable.

\paragraph{Token transform.}
\begin{equation}
  y_{i,j} \;=\; \phi(x_{i,j}) \;=\; \sigma\!\big(W_1 x_{i,j} + b_1\big),
  \qquad W_1\in\mathbb{R}^{m\times d},
  \label{eq:mlp}
\end{equation}
with $\sigma$ a pointwise nonlinearity (GELU below).

\paragraph{Aggregators.} Write $\alpha_{i,j}\ge 0$, $\sum_{j}\alpha_{i,j}=1$ for a
convex pooling weight. The four rules compared are
%
\begin{align}
  \text{last token:}\quad
    &\alpha_{i,j}=\mathbb{1}[\,j=N_i\,], \label{eq:agg-last}\\[2pt]
  \text{mean:}\quad
    &\alpha_{i,j}=1/N_i, \label{eq:agg-mean}\\[2pt]
  \text{softmax attention:}\quad
    &\alpha_{i,j}=\frac{\exp(q^\top y_{i,j}/\sqrt{m})}
                       {\sum_{j'=1}^{N_i}\exp(q^\top y_{i,j'}/\sqrt{m})},
    \label{eq:agg-softmax}\\[2pt]
  \text{exponential moving average:}\quad
    &\alpha_{i,j}=\frac{\lambda^{\,N_i-j}}{\sum_{j'=1}^{N_i}\lambda^{\,N_i-j'}},
     \quad \lambda\in(0,1). \label{eq:agg-ema}
\end{align}
%
For \eqref{eq:agg-last}--\eqref{eq:agg-ema} the logit is $w^\top\!\big(\sum_j
\alpha_{i,j}y_{i,j}\big)+b$. \textbf{MultiMax} instead applies a hard maximum
per head, which is \emph{not} of this convex-combination form:
%
\begin{equation}
  \boxed{\;
  a_h(X_i) \;=\; \max_{1\le j\le N_i}\; v_h^\top y_{i,j},
  \qquad
  \mathrm{logit}_i \;=\; \sum_{h=1}^{H} a_h(X_i) \;+\; b,
  \;}
  \label{eq:multimax}
\end{equation}
with heads $v_h\in\mathbb{R}^{m}$, $h=1,\dots,H$. Predictions are obtained through a
bounded link,
%
\begin{equation}
  \hat p_i \;=\; \varsigma\!\big(\Pi_{[-c,\,c]}(\mathrm{logit}_i)\big),
  \qquad \varsigma(z)=\big(1+e^{-z}\big)^{-1},\quad c=10,
  \label{eq:clamped-sigmoid}
\end{equation}
where $\Pi_{[-c,c]}$ is the projection onto $[-c,c]$.

\begin{remark}[Gradient of the clamp]
\label{rem:clamp}
$\Pi_{[-c,c]}$ has $\partial\Pi/\partial z=\mathbb{1}[\,|z|<c\,]$, which is
\emph{identically zero} on the saturated region. Implemented naively, the guard
introduced to avoid vanishing gradients produces them. We therefore use the
straight-through estimator
%
\begin{equation}
  \widetilde\Pi(z)\;=\;z+\operatorname{sg}\!\big(\Pi_{[-c,c]}(z)-z\big),
  \qquad
  \frac{\partial\widetilde\Pi}{\partial z}\equiv 1,
  \label{eq:straight-through}
\end{equation}
with $\operatorname{sg}$ the stop-gradient. The forward value is exactly the clamp in
\eqref{eq:clamped-sigmoid}; the backward pass is the identity.
\end{remark}


\section{Context Dilution}
\label{sec:dilution}

We formalise why averaging aggregators lose a localised signal as context grows, and why
\eqref{eq:multimax} does not.

\begin{assumption}[Localised attack in benign context]
\label{asm:needle}
A sequence of length $N$ contains an attack segment $\mathcal{A}\subset\{1,\dots,N\}$
with $|\mathcal{A}|=k$ fixed and independent of $N$. There exist constants
$\mu_{\mathrm{a}}>\mu_{\mathrm{b}}$ such that the head-projected features satisfy
\begin{equation}
  v_h^\top y_{j} \;=\;
  \begin{cases}
    \mu_{\mathrm{a}} + \varepsilon_j, & j\in\mathcal{A},\\
    \mu_{\mathrm{b}} + \varepsilon_j, & j\notin\mathcal{A},
  \end{cases}
  \label{eq:needle-model}
\end{equation}
with $\varepsilon_j$ i.i.d., zero-mean, sub-Gaussian with variance proxy $\tau^2$.
Write $\Delta:=\mu_{\mathrm{a}}-\mu_{\mathrm{b}}>0$ for the signal gap.
\end{assumption}

\subsection{Mean aggregation}

\begin{proposition}[Mean pooling dilutes as $\Theta(1/N)$]
\label{prop:mean}
Under Assumption~\ref{asm:needle}, the mean-pooled head score satisfies
\begin{equation}
  \mathbb{E}\Big[\tfrac{1}{N}\textstyle\sum_{j=1}^{N} v_h^\top y_j\Big]
    \;-\;\mu_{\mathrm{b}}
  \;=\; \frac{k}{N}\,\Delta \;=\; \Theta\!\big(N^{-1}\big),
  \label{eq:mean-dilution}
\end{equation}
while the pooled fluctuation has standard deviation $\tau/\sqrt{N}$. The resulting
signal-to-noise ratio is
\begin{equation}
  \mathrm{SNR}_{\mathrm{mean}}(N)
  \;=\;\frac{(k/N)\Delta}{\tau/\sqrt{N}}
  \;=\;\frac{k\Delta}{\tau}\,N^{-1/2}
  \;\xrightarrow[N\to\infty]{}\;0 .
  \label{eq:mean-snr}
\end{equation}
\end{proposition}

\begin{proof}
Linearity of expectation gives the mean shift $\big(k\mu_{\mathrm{a}}+(N-k)\mu_{\mathrm{b}}\big)/N-\mu_{\mathrm{b}}=k\Delta/N$. Independence gives variance $\tau^2/N$. Dividing yields \eqref{eq:mean-snr}.
\end{proof}

\begin{corollary}[Detection collapse]
\label{cor:mean-collapse}
For any fixed decision threshold calibrated at length $N_0$, the mean-pooled detector's
power at length $N\gg N_0$ decays as $N^{-1/2}$; equivalently, holding a target
true-positive rate fixed requires $\Delta$ to grow like $\sqrt{N}$.
\end{corollary}

\subsection{Softmax aggregation}

\begin{proposition}[Softmax attention weight on the attack is $\mathcal{O}(1/N)$ for
bounded logit gaps]
\label{prop:softmax}
Let the attention logits be $s_j=q^\top y_j/\sqrt{m}$ with
$s_j\le \bar{s}$ for $j\in\mathcal{A}$ and $s_j\ge \underline{s}$ for
$j\notin\mathcal{A}$, and set $\gamma:=\bar{s}-\underline{s}<\infty$. Then the total
attention mass on the attack segment obeys
\begin{equation}
  \alpha_{\mathcal{A}}
  \;:=\;\sum_{j\in\mathcal{A}}\alpha_j
  \;\le\; \frac{k\,e^{\gamma}}{k\,e^{\gamma}+(N-k)}
  \;=\;\mathcal{O}\!\big(N^{-1}\big),
  \qquad N\to\infty .
  \label{eq:softmax-mass}
\end{equation}
\end{proposition}

\begin{proof}
Write $Z=\sum_{j}e^{s_j}$. Bounding the numerator above and the benign part of the
denominator below,
\[
  \alpha_{\mathcal{A}}
  =\frac{\sum_{j\in\mathcal{A}}e^{s_j}}{Z}
  \le\frac{k\,e^{\bar{s}}}{k\,e^{\bar{s}}+(N-k)e^{\underline{s}}}
  =\frac{k\,e^{\gamma}}{k\,e^{\gamma}+(N-k)} ,
\]
after dividing through by $e^{\underline{s}}$. The right-hand side is
$k e^{\gamma}/N+\mathcal{O}(N^{-2})$.
\end{proof}

\begin{remark}[Why the gap is bounded in practice]
\label{rem:gamma}
Proposition~\ref{prop:softmax} is vacuous if $\gamma$ may grow with $N$. It does not:
$\gamma$ is set by the trained query $q$ and the feature geometry, both fixed at
deployment, while $N$ is chosen by the adversary. Escaping dilution would require
$\gamma\gtrsim\log N$, i.e.\ attention logits growing without bound --- which the
$1/\sqrt{m}$ scaling in \eqref{eq:agg-softmax} is specifically designed to prevent.
This is the formal content of the empirical observation that softmax pooling
\emph{flattens} as context grows.
\end{remark}

\subsection{MultiMax aggregation}

\begin{theorem}[Padding invariance of the hard maximum]
\label{thm:multimax}
Let $X\in\mathbb{R}^{N\times d}$ contain an attack segment $\mathcal{A}$ and let
$X'=\mathrm{pad}(X,P)$ append $P$ arbitrary benign tokens. Then for every head $h$,
\begin{equation}
  a_h(X')\;=\;\max\Big\{\,a_h(X),\;\max_{j\in\text{pad}} v_h^\top y_j\,\Big\}
  \;\ge\;a_h(X).
  \label{eq:multimax-monotone}
\end{equation}
In particular, if the padded tokens satisfy
$v_h^\top y_j\le a_h(X)$ --- which holds whenever the padding is benign in the sense of
Assumption~\ref{asm:needle} and $\Delta>0$ dominates the benign extremes --- then
\begin{equation}
  \boxed{\;a_h(X')\;=\;a_h(X)\;=\;\Omega(1),\quad\text{independent of }P.\;}
  \label{eq:multimax-invariance}
\end{equation}
\end{theorem}

\begin{proof}
The maximum over a union of index sets is the maximum of the sub-maxima, giving
\eqref{eq:multimax-monotone}; monotonicity is immediate since
$\mathcal{A}\subseteq\{1,\dots,N\}\subseteq\{1,\dots,N+P\}$. Under the stated condition
the second term does not exceed the first, so the maximum is attained inside the
original segment and its value is unchanged.
\end{proof}

\begin{corollary}[Asymptotic separation of the aggregators]
\label{cor:separation}
Under Assumption~\ref{asm:needle}, as $N\to\infty$ with $k$ fixed,
\begin{equation}
  \frac{\mathrm{signal}_{\mathrm{multimax}}(N)}{\mathrm{signal}_{\mathrm{mean}}(N)}
  \;=\;\Theta(N),
  \qquad
  \frac{\mathrm{signal}_{\mathrm{multimax}}(N)}{\mathrm{signal}_{\mathrm{softmax}}(N)}
  \;=\;\Omega(N).
  \label{eq:separation}
\end{equation}
\end{corollary}

\begin{remark}[The cost of the maximum: benign extremes]
\label{rem:fpr}
Theorem~\ref{thm:multimax} is one-sided. The same monotonicity that protects recall
also means the benign maximum grows with $N$: for sub-Gaussian $\varepsilon_j$,
$\mathbb{E}\big[\max_{j\le N}\varepsilon_j\big]=\Theta\big(\tau\sqrt{2\log N}\big)$.
A MultiMax probe therefore has a false-positive rate that drifts upward with context
length unless its threshold is recalibrated as
\begin{equation}
  t(N)\;=\;t_0+\tau\sqrt{2\log N},
  \label{eq:threshold-drift}
\end{equation}
which is a $\sqrt{\log N}$ correction rather than the $\sqrt{N}$ correction
Corollary~\ref{cor:mean-collapse} demands of mean pooling. This is the honest statement
of the trade: MultiMax converts a polynomial degradation into a logarithmic one; it does
not eliminate length dependence.
\end{remark}

\begin{remark}[Gradient sparsity: the optimisation cost of the hard maximum]
\label{rem:gradient-sparsity}
Theorem~\ref{thm:multimax} concerns the \emph{forward} map only. In training, the
subgradient of $a_h$ is supported on a single position,
\begin{equation}
  \frac{\partial a_h}{\partial y_{j}}
  \;=\;v_h\,\mathbb{1}\!\left[\,j=\argmax_{j'} v_h^\top y_{j'}\,\right],
  \label{eq:hardmax-grad}
\end{equation}
so each head receives learning signal from exactly one token per sequence, whereas
softmax pooling \eqref{eq:agg-softmax} spreads gradient over all $N$ positions with
weights $\alpha_j$. The effective sample size per gradient step is therefore $H$ tokens
rather than $N$. Near the detection threshold --- where the argmax is not yet reliably
inside $\mathcal{A}$ --- this can prevent the probe from learning the concept at all,
even at short training lengths where dilution is irrelevant.

This predicts a non-monotone comparison: MultiMax should \emph{lose} to softmax pooling
in a band of weak signal strengths, and match it once the signal is strong enough for the
argmax to lock on. Measured (\S\,benchmark, $k=4$, $N=16{,}384$, recall at $1\%$ FPR,
reported as long-context / training-length):
%
\begin{center}
\begin{tabular}{lccc}
\hline
strength & MultiMax & softmax & mean \\
\hline
$0.10$ & $0.00\,/\,0.00$ & $\mathbf{0.62\,/\,1.00}$ & $0.00\,/\,0.47$\\
$0.15$ & $1.00\,/\,1.00$ & $1.00\,/\,1.00$ & $0.00\,/\,0.95$\\
$0.50$ & $1.00\,/\,1.00$ & $1.00\,/\,1.00$ & $0.40\,/\,1.00$\\
\hline
\end{tabular}
\end{center}
%
At strength $0.10$ MultiMax fails to learn at its own training length while softmax
attains $1.00$ there and retains $0.62$ under a $32\times$ context extension. Both cross
the $99\%$ threshold at strength $0.15$; mean pooling never does.
\end{remark}

\begin{remark}[Smooth-max annealing removes the optimisation failure]
\label{rem:annealing}
The sparsity in \eqref{eq:hardmax-grad} is a property of the \emph{training} objective,
not of the deployed map, so it can be annealed away. Replace the hard maximum by the
normalised Boltzmann operator
\begin{equation}
  \mathrm{smax}_\tau(s)
  \;=\;\tau\log\!\Big(\tfrac{1}{N}\textstyle\sum_{j=1}^{N} e^{s_j/\tau}\Big),
  \qquad
  \mathrm{smax}_\tau(s)\;\xrightarrow[\tau\to 0]{}\;\max_j s_j,
  \quad
  \mathrm{smax}_\tau(s)\;\xrightarrow[\tau\to\infty]{}\;\tfrac1N\textstyle\sum_j s_j,
  \label{eq:boltzmann}
\end{equation}
whose gradient $\partial\,\mathrm{smax}_\tau/\partial s_j=\mathrm{softmax}(s/\tau)_j$ is
dense for $\tau>0$. Annealing $\tau\to 0$ over training recovers the exact hard max at
deployment, so Theorem~\ref{thm:multimax} still applies to the shipped model.

\begin{center}
\begin{tabular}{lccc}
\hline
strength & MultiMax (no anneal) & MultiMax (annealed) & softmax \\
\hline
$0.10$ & $0.00\,/\,0.00$ & $\mathbf{1.00\,/\,1.00}$ & $0.68\,/\,1.00$\\
$0.15$ & $1.00\,/\,1.00$ & $1.00\,/\,1.00$ & $1.00\,/\,1.00$\\
$0.50$ & $1.00\,/\,1.00$ & $1.00\,/\,1.00$ & $1.00\,/\,1.00$\\
\hline
\end{tabular}
\end{center}

\textbf{The $1/N$ normalisation in \eqref{eq:boltzmann} is not cosmetic.} Plain
$\mathrm{LSE}_\tau$ adds $H\tau\log N$ to the logit --- measured at $+42$ for $\tau=1$,
$N=512$ and $+68$ at $N=16{,}384$. That offset is both large and \emph{length-dependent},
so an annealing probe must relearn its bias each epoch and meets a different offset at
deployment length than at training length. The observable symptom was a recall curve
non-monotone in signal strength ($1.00$ at $s=0.15$, $0.04$ at $s=0.50$).
\end{remark}

\begin{remark}[Distributed attacks: the structural limit of the maximum]
\label{rem:distributed}
Theorem~\ref{thm:multimax} assumes the attack mass is concentrated. An adversary who
holds the total budget $S$ fixed but spreads it over $m$ non-contiguous tokens presents
per-token magnitude $S/m$, so
\begin{equation}
  \max_j v_h^\top y_j \;=\;\Theta(S/m),
  \qquad\text{while}\qquad
  \tfrac1N\textstyle\sum_j v_h^\top y_j \;=\;\Theta(S/N)\ \ \text{is unchanged in } m .
  \label{eq:distributed}
\end{equation}
The max degrades in $m$ while the mean does not: this is the exact dual of
Proposition~\ref{prop:mean}, and it is the structural price of the hard reduction.

Measured at $S=2.0$, $N=16{,}384$, recall at $1\%$ FPR:
\begin{center}
\begin{tabular}{lcccc}
\hline
$m$ & per-token $S/m$ & MultiMax & softmax & mean \\
\hline
$1$   & $2.0000$ & $1.00$ & $1.00$ & $0.14$\\
$16$  & $0.1250$ & $1.00$ & $1.00$ & $0.14$\\
$64$  & $0.0312$ & $0.48$ & $0.44$ & $0.10$\\
$256$ & $0.0078$ & $0.02$ & $0.00$ & $0.06$\\
\hline
\end{tabular}
\end{center}

Failure boundary (lowest $m$ with recall $<0.80$): MultiMax $m=64$, softmax $m=64$,
mean $m=1$.

Two cautions. First, the predicted MultiMax-specific vulnerability is \emph{not} isolated
here: both pooled aggregators fail at the same $m$, so at this budget the distributed
attack degrades the query-based reduction as much as the hard one. A preliminary run at
$N=8{,}192$ suggested softmax survived to $m=256$; that did \emph{not} replicate at
$N=16{,}384$ and the earlier figure should be disregarded. Second, the boundary hides a
real difference in \emph{ranking} quality: at $m=256$ both have recall $\approx 0$ at a
$1\%$ FPR threshold, but AUROC is $0.523$ for MultiMax against $0.716$ for softmax, so
softmax retains usable ordering after the thresholded detector has failed. Under a
distributed adversary the hard max loses its signal more completely, even where the
headline recall numbers coincide.
\end{remark}

\begin{remark}[Scope of the dilution advantage]
\label{rem:scope}
Combining Prop.~\ref{prop:mean}, Prop.~\ref{prop:softmax},
Remark~\ref{rem:annealing} and Remark~\ref{rem:distributed}, the honest statement is
three-part.

\emph{(i) Against mean aggregation, MultiMax wins unconditionally on concentrated
attacks.} The $\Theta(1/N)$ decay of Prop.~\ref{prop:mean} has no free parameter to
absorb it, and mean pooling fails at every measured signal strength.

\emph{(ii) Against single-query softmax pooling, the advantage is conditional and
narrow.} $\gamma$ in \eqref{eq:softmax-mass} is a \emph{trained} quantity, and a
well-separated needle drives $e^{\gamma}\!\sim\!N/k$, cancelling the dilution exactly.
Remark~\ref{rem:gamma} argues $\gamma$ cannot grow without bound; the measurements show
that at $N=16{,}384$ it does not need to. Post-annealing, MultiMax leads softmax only in
the weak-signal regime ($1.00$ vs $0.68$ at $s=0.10$) and ties elsewhere. Claims of
superiority should be stated as conditional on $k/N$ and the achievable logit gap.

\emph{(iii) Against a distributed adversary, both pooled aggregators fail together}
(Remark~\ref{rem:distributed}): the recall boundary is $m=64$ for MultiMax and for
softmax alike. The hard max nonetheless loses its ranking signal more completely
(AUROC $0.52$ vs $0.72$ at $m=256$), so the distributed attack is a limitation of the
whole approach at this budget rather than a MultiMax-specific weakness --- with MultiMax
degrading less gracefully once past the boundary.

The unconditional advantage is therefore the systems one --- $\Theta(1)$ activation
memory in $N$ by Prop.~\ref{prop:complexity}, measured as $35$\,MiB of overhead at
$N=131{,}072$ against $652$\,MiB for softmax pooling --- and not a statistical one.
\end{remark}

\subsection{Complexity}

\begin{proposition}[Memory and time]
\label{prop:complexity}
For batch size $1$, sequence length $N$, width $m$, heads $H$, and chunk size $C$:
\begin{equation}
  \begin{aligned}
    \text{MultiMax (chunked):}\quad &\Theta(Nmd)\ \text{time},\quad
      \Theta(Cm+H)\ \text{activation memory},\\
    \text{softmax attention (single query):}\quad &\Theta(Nmd)\ \text{time},\quad
      \Theta(Nm)\ \text{activation memory},\\
    \text{self-attention:}\quad &\Theta(N^2m)\ \text{time},\quad
      \Theta(N^2)\ \text{activation memory}.
  \end{aligned}
  \label{eq:complexity}
\end{equation}
\end{proposition}

\begin{remark}[Where the quadratic cost actually lives, and what the systems win is]
\label{rem:complexity-scope}
Single-query attention pooling is $\Theta(N)$ in memory, not $\Theta(N^2)$: its score
tensor has shape $(1,N)$. The quadratic term appears only when the probe performs
self-attention \emph{across} the sequence.

Against single-query pooling, then, MultiMax offers no asymptotic memory separation ---
both are $\Theta(N)$ if the input is counted, and the input is unavoidable. The win is a
large constant: streaming the reduction keeps the aggregator's \emph{own} overhead at
$\Theta(1)$ in $N$. Measured overhead above the input tensor, $d=2048$, batch $1$:
%
\begin{center}
\begin{tabular}{lccccc}
\hline
$N$ & $128$ & $1{,}024$ & $8{,}192$ & $32{,}768$ & $131{,}072$\\
\hline
MultiMax        & $11.4$ & $13.1$ & $19.1$ & $19.1$ & $\mathbf{19.1}$ MiB\\
softmax pooling & $11.8$ & $18.1$ & $63.2$ & $171.4$ & $\mathbf{652.1}$ MiB\\
self-attention  & $14.5$ & $28.0$ & $685.6$ & $10381.6$ & OOM\\
\hline
\end{tabular}
\end{center}
%
with log--log exponents $\alpha=0.000$, $0.842$ and $1.960$ respectively (fitted on
$N\ge 8192$). MultiMax overhead is flat across a $16\times$ length increase, and
self-attention is $492\times$ slower at $N=32{,}768$ before exhausting memory entirely.

Note that this is the \emph{unconditional} advantage. By Remark~\ref{rem:scope} the
statistical advantage over softmax pooling is conditional and narrow, so the case for
the hard reduction rests here rather than on Theorem~\ref{thm:multimax} alone.
\end{remark}


\section{Cost--Accuracy Optimisation for Uncertainty-Gated Cascades}
\label{sec:cascade}

\subsection{Routing rule and system cost}

Let $p_i=\varsigma(A\,z_i+B)$ be the calibrated probe probability, where $z_i$ is the raw
logit and $(A,B)$ are Platt parameters (\S\ref{sec:calibration}). Fix $\delta>0$ and
define the ambiguous band
\begin{equation}
  \mathcal{U}_\delta \;=\;\big\{\,i:\ |p_i-\tfrac12|<\delta\,\big\},
  \qquad
  u(\delta)\;=\;\Pr\big[\,|p-\tfrac12|<\delta\,\big].
  \label{eq:band}
\end{equation}
The cascade returns the probe's decision outside $\mathcal{U}_\delta$ and defers to the
LLM inside it. With $c_{\mathrm{p}}$ the per-request probe cost and
$c_{\mathrm{L}}$ the per-request LLM cost,
\begin{equation}
  \boxed{\;
  C_{\mathrm{total}}(\delta)\;=\;c_{\mathrm{p}}\;+\;u(\delta)\,c_{\mathrm{L}}\;}
  \label{eq:cost}
\end{equation}
per request, against the LLM-only baseline $C_{\mathrm{base}}=c_{\mathrm{L}}$. The
relative saving is
\begin{equation}
  S(\delta)\;=\;1-\frac{C_{\mathrm{total}}(\delta)}{C_{\mathrm{base}}}
  \;=\;\big(1-u(\delta)\big)-\frac{c_{\mathrm{p}}}{c_{\mathrm{L}}} .
  \label{eq:saving}
\end{equation}
Since $c_{\mathrm{p}}/c_{\mathrm{L}}\ll 1$ in the regime of interest, the saving is
governed almost entirely by the escalation rate $u(\delta)$.

\subsection{Accuracy of the cascade}

Let $\mathrm{Acc}_{\mathrm{p}}(\mathcal{S})$ and $\mathrm{Acc}_{\mathrm{L}}(\mathcal{S})$
denote probe and LLM accuracy on a region $\mathcal{S}$. Then
\begin{equation}
  \mathrm{Acc}(\delta)
  \;=\;\big(1-u(\delta)\big)\,\mathrm{Acc}_{\mathrm{p}}\!\big(\mathcal{U}_\delta^{c}\big)
      \;+\;u(\delta)\,\mathrm{Acc}_{\mathrm{L}}\!\big(\mathcal{U}_\delta\big).
  \label{eq:acc}
\end{equation}
Differentiating \eqref{eq:acc} and \eqref{eq:cost} with respect to $\delta$ and
eliminating $\mathrm{d}\delta$ gives the exchange rate between accuracy and money:
\begin{equation}
  \frac{\mathrm{d}\,\mathrm{Acc}}{\mathrm{d}\,C_{\mathrm{total}}}
  \;=\;\frac{1}{c_{\mathrm{L}}}\Big[
      \mathrm{Acc}_{\mathrm{L}}\!\big(\partial\mathcal{U}_\delta\big)
      -\mathrm{Acc}_{\mathrm{p}}\!\big(\partial\mathcal{U}_\delta\big)\Big],
  \label{eq:exchange}
\end{equation}
evaluated on the boundary shell $\partial\mathcal{U}_\delta=\{i:|p_i-\frac12|=\delta\}$.

\begin{proposition}[Optimal band width]
\label{prop:optimal-delta}
For a budget $\Lambda$ (dollars per unit accuracy), the Lagrangian
$\mathcal{L}(\delta)=\mathrm{Acc}(\delta)-\Lambda\,C_{\mathrm{total}}(\delta)$ is
stationary where
\begin{equation}
  \mathrm{Acc}_{\mathrm{L}}\!\big(\partial\mathcal{U}_\delta\big)
  -\mathrm{Acc}_{\mathrm{p}}\!\big(\partial\mathcal{U}_\delta\big)
  \;=\;\Lambda\,c_{\mathrm{L}} .
  \label{eq:stationarity}
\end{equation}
Escalation is worthwhile exactly while the LLM's marginal advantage on the shell exceeds
its marginal price. Because $\mathrm{Acc}_{\mathrm{p}}\to\frac12$ as $\delta\to0$ (the
probe is maximally uncertain at the boundary) and the advantage shrinks monotonically as
$\delta$ grows, \eqref{eq:stationarity} has a unique root under a unimodal score density,
and the Pareto frontier exhibits a knee at that $\delta^\ast$.
\end{proposition}

\subsection{Calibration is a precondition, not a refinement}
\label{sec:calibration}

The routing rule \eqref{eq:band} is a statement about $p$, so it is meaningful only if
$p$ is calibrated. Two failure modes bracket the design.

\begin{remark}[Degenerate gating]
\label{rem:degenerate}
If the score distribution is over-dispersed, $u(\delta)\to0$ and the cascade silently
becomes probe-only at full cost saving and probe accuracy. If it is under-dispersed,
$u(\delta)\to1$ and the cascade becomes LLM-only, paying the baseline while appearing to
be gated. Both are invisible in $C_{\mathrm{total}}$ alone; both are detected by
reporting $u(\delta)$.
\end{remark}

\begin{remark}[Why temperature scaling is insufficient for MultiMax]
\label{rem:platt}
The MultiMax logit \eqref{eq:multimax} is a sum of $H$ \emph{maxima}. By
Remark~\ref{rem:fpr}, $\mathbb{E}[a_h]=\Theta(\sqrt{2\log N})$ even under the null, so
the logit carries an $N$-dependent positive offset. Temperature scaling
$p=\varsigma(z/T)$ is a pure rescaling and fixes $\varsigma^{-1}(p)=0$ at $z=0$; it
cannot move a mis-centred decision boundary. The two-parameter Platt map
\begin{equation}
  p \;=\; \varsigma\!\big(A\,z+B\big),
  \qquad
  (A,B)\;=\;\argmin_{a,b}\;-\!\!\sum_i\Big[y_i\log\varsigma(az_i+b)+(1-y_i)\log\big(1-\varsigma(az_i+b)\big)\Big],
  \label{eq:platt}
\end{equation}
absorbs both the scale and the offset, and is the minimal correction that makes
$|p-\frac12|$ an interpretable quantity for \eqref{eq:band}.
\end{remark}

\subsection{Aggregate cost model}

For $R$ requests with mean prompt length $\bar{n}$ tokens, unit prices $\pi_{\mathrm{in}}$
and $\pi_{\mathrm{out}}$ per token and $n_{\mathrm{out}}$ output tokens per verdict,
\begin{equation}
  C_{\mathrm{total}}(\delta;R)
  \;=\;R\Big[\underbrace{\kappa\,\bar{n}}_{\text{probe}}
    \;+\;u(\delta)\big(\pi_{\mathrm{in}}\bar{n}+\pi_{\mathrm{out}}n_{\mathrm{out}}\big)\Big],
  \label{eq:aggregate-cost}
\end{equation}
with $\kappa$ the amortised accelerator cost per token. Tokens routed to the LLM are
$T(\delta)=R\,u(\delta)\,\bar{n}$, so token savings relative to the baseline are
\begin{equation}
  1-\frac{T(\delta)}{R\,\bar{n}}\;=\;1-u(\delta),
  \label{eq:token-saving}
\end{equation}
i.e.\ the escalation rate is simultaneously the token-saving deficit and, up to the
negligible $c_{\mathrm{p}}/c_{\mathrm{L}}$ term in \eqref{eq:saving}, the monetary one.

% =====================================================================================
% End of math_formulation.tex
% =====================================================================================
